\section{Second-order formulas}
\label{sec:second-order-logic-01}

\begin{definition}[Second-order signature]
A second-order signature \(\Lang_2\) contains the constant, function, and
relation symbols of a first-order signature.  Its logical alphabet also
contains, for every \(n\geq1\), a countably infinite set
\(\mathsf{RelVar}_n\) of \(n\)-ary relation variables.  Relation variables are
logical variables; nonlogical relation symbols remain members of \(\Lang_2\).
\end{definition}

First-order terms are generated as before.  If \(X\in\mathsf{RelVar}_n\) and
\(t_1,\ldots,t_n\) are terms, then \(X(t_1,\ldots,t_n)\) is atomic.  Formulas
are closed under the propositional connectives, first-order quantifiers
\(\forall x,\exists x\), and second-order quantifiers
\(\forall X^n,\exists X^n\).  The superscript records arity and may be omitted
only when the arity is fixed by context.

\begin{definition}[Second-order sentence]
A second-order sentence is a second-order formula having no free first-order
variable occurrence and no free relation-variable occurrence.
\end{definition}

Term substitution for first-order variables and relation-formula substitution
for relation variables are different operations.  For example, substituting a
formula \(\varphi(x_1,\ldots,x_n)\) for a free occurrence of \(X^n\) requires
renaming every binder that could capture a variable free in \(\varphi\).

